\ifdefined\gkver\else\def\gkver{student}\fi
\documentclass[\gkver]{gaokaozhenti}
\begin{document}

\chapter{2005年上海理综}
%% paperType: 理综物理部分

\begin{enumerate}
\item
%% number: 1
%% paperName: 2005年上海理综第 1 题 3 分
%% typeId: 1
%% type: 单选题
%% chapter: 运动和力
%% point: 牛顿第一定律
%% method: 无
%% score: 3
%% degree: 850
%% duplicateId: 0
%% body:
下列词语所隐含的物理现象，可用牛顿理论中力的概念和规律来解释的是\xzanswer{C}

①滴水穿石　②铁杵成针　③水中望月　④逆风扬帆　⑤光阴飞逝

\fourchoices[answer=C]
{①②③}
{③④⑤}
{①②④}
{①③⑤}

%% answer: C
%% memo:
\memoanswer{
本题原卷及材料中未提供详解，待补充。
}

\item
%% number: 5
%% paperName: 2005年上海理综第 5 题 3 分
%% typeId: 1
%% type: 单选题
%% chapter: 其他
%% point: 物理思想
%% method: 类比
%% score: 3
%% degree: 650
%% duplicateId: 0
%% body:
关于原子结构，汤姆生提出葡萄干蛋糕模型、卢瑟福提出行星模型$\ldots\ldots$，都采用了类比推理的方法。下列事实中，主要采用类比推理的是\xzanswer{C}

\fourchoices[answer=C]
{人们为便于研究物体的运动而建立的质点模型}
{伽利略从教堂吊灯的摆动中发现摆的等时性规律}
{库仑根据牛顿的万有引力定律提出库仑定律}
{托马斯$\cdot$杨通过双缝干涉实验证实光是一种波}

%% answer: C
%% memo:
\memoanswer{
本题原卷及材料中未提供详解，待补充。
}

\item
%% number: 31
%% paperName: 2005年上海理综第 31 题 3 分
%% typeId: 3
%% type: 填空题
%% chapter: 运动和力
%% point: 牛顿定律局限性
%% method: 无
%% score: 3
%% degree: 850
%% duplicateId: 0
%% body:
右图是 2005 年“世界物理年”的徽标。观察这个徽标，提出一个问题，并写在下面的横线上。（例：为什么把 2005 年确定为世界物理年？）\tkanswer[8cm]{徽标中的沙漏（上、下圆圈、中间的叉、整体的锥形等）有什么样的物理含义？}。

\onepicture[w=3.2cm, align=right]{figs/31.png}

%% answer: 
\jdanswer{
徽标中的沙漏（上、下圆圈、中间的叉、整体的锥形等）有什么样的物理含义？
}
%% memo:
\memoanswer{
本题原卷及材料中未提供详解，待补充。
}

\item
%% number: 33
%% paperName: 2005年上海理综第 33 题 4 分
%% typeId: 3
%% type: 填空题
%% chapter: 运动和力
%% point: 牛顿定律局限性
%% method: 无
%% score: 4
%% degree: 650
%% duplicateId: 0
%% body:
“世界物理年”决议的作出与爱因斯坦的相对论时空观有关。根据爱因斯坦的理论，一把米尺，在它与观察者有不同相对速度情况下，米尺长度是不同的，它们之间的关系见右图。由此可知，当米尺与观察者的相对速度达到 $0.8c$（$c$ 为光速）时，米尺长度大约是 \tkanswer[2cm]{0.6} 米。在日常生活中，我们无法察觉到米尺长度变化的现象，是因为观察者相对于米尺的运动速度 \tkanswer[5cm]{远远小于光速或速度很小}。

\onepicture[w=7.6cm, align=right]{figs/33.png}

%% answer: 
\jdanswer{
0.6；远远小于光速或速度很小
}
%% memo:
\memoanswer{
本题原卷及材料中未提供详解，待补充。
}

\end{enumerate}

\end{document}
